% ==============================
% Chapter 2 – Problem Definition
% ==============================

\chapter{Problem Definition}
% --------------------------------
% 2.1 Problem Statement
% --------------------------------

\section{Problem Statement}

The primary problem this project addresses is \textbf{hardware fragmentation} and \textbf{vendor lock-in} in the quantum computing domain. Currently, quantum hardware providers use different underlying technologies (e.g., Superconducting Qubits vs. Trapped Ions), which result in incompatible instruction sets and connectivity graphs. A developer writing code for one platform often cannot run it on another without significant manual refactoring. Furthermore, optimizing a circuit for a specific device's topology is a complex, error-prone task that requires deep knowledge of the hardware. There is a lack of lightweight, educational, and open-source runtimes that demonstrate how to abstract these complexities effectively. This project aims to solve these interoperability challenges by automating the translation process, thereby reducing development time, eliminating vendor dependencies, and lowering the barrier to entry for quantum computing education and research.



% --------------------------------
% 2.2 Proposed Solution Overview
% --------------------------------

\section{Proposed Solution Overview}

The proposed system solves the hardware fragmentation problem by introducing a middleware layer—the Quantum Virtual Machine (QVM). Instead of targeting a specific device, developers write code for the QVM, which follows a three-stage pipeline: (1) \textbf{Ingestion:} The system accepts a high-level circuit (OpenQASM 3.0, 2.0, or JSON) and converts it into a standardized Intermediate Representation (IR). (2) \textbf{Transpilation:} The core engine analyzes the target backend's constraints (native gates and topology), inserts SWAP gates where necessary, and decomposes complex gates into the target's primitive basis gates. (3) \textbf{Execution:} The system either executes the circuit on its internal hybrid simulator (Statevector or MPS) for verification or outputs the compiled circuit for a specific hardware target. This approach ensures that the logic of the algorithm remains preserved while the implementation details are automatically handled by the software, achieving true "Write Once, Run Anywhere" (WORA) functionality.



% --------------------------------
% 2.3 Objectives of the Proposed System
% --------------------------------

\section{Objectives of the Proposed System}

The following business objectives guide the development and evaluation of the Quantum Virtual Machine:

\vspace{0.3cm}

\textbf{Business Objectives (BO):}

\begin{itemize}
    \item \textbf{BO-1:} Develop a hardware-agnostic quantum execution runtime that enables developers to write quantum algorithms once and transpile them for different qubit topologies (Linear, Grid) without rewriting code.
    
    \item \textbf{BO-2:} Provide educational transparency into the transpilation process, showing users how logical qubits are mapped to physical ones and where SWAP gates are inserted to respect hardware constraints.
    
    \item \textbf{BO-3:} Implement a lightweight, standalone runtime suitable for local testing of small-scale quantum algorithms (up to 10-20 qubits) without requiring paid cloud subscriptions or enterprise dependencies.
    
    \item \textbf{BO-4:} Ensure interoperability by supporting OpenQASM 2.0 and 3.0 standards, making compiled circuits compatible with IBM quantum processors and other QASM-compliant platforms.
    
    \item \textbf{BO-5:} Automate the complex tasks of gate decomposition, qubit mapping, and circuit optimization through a robust transpilation pipeline, reducing manual effort and eliminating the need for deep hardware-specific knowledge.
\end{itemize}



% --------------------------------
% 2.4 Scope of the System
% --------------------------------

\section{Scope of the System}

The scope of this project is to develop a desktop-based Quantum Virtual Machine (QVM) using Python. The system will support the creation, compilation, and simulation of quantum circuits with a maximum capacity of 10-20 qubits depending on the simulation backend. The core functionality centers on a "Transpiler Pipeline" that accepts a high-level circuit definition (OpenQASM 3.0, 2.0, or JSON) and adapts it to specific architectural constraints (e.g., restricted connectivity). The system will support both Statevector simulation (exact linear algebra for up to 12 qubits) and MPS (Matrix Product States) simulation for low-entanglement circuits scaling to 20+ qubits. The output will include the final state vector, measurement probabilities, classical memory states, and visualized circuit diagrams. The project will \textbf{not} include pulse-level control, quantum error correction (QEC), or direct integration with real hardware APIs beyond OpenQASM string generation. The target users are quantum computing students, researchers, and educators requiring a lightweight, transparent, and educational quantum runtime environment.



% --------------------------------
% 2.5 System Architecture and Modules
% --------------------------------

\section{System Architecture and Modules}

The QVM follows a strict \textbf{Pipeline Architecture} pattern, ensuring clear separation of concerns and modular extensibility. The system processes quantum programs through six sequential stages: (1) \textbf{Input Layer:} Accepts OpenQASM 3.0, 2.0, or JSON gate lists. (2) \textbf{Parser (Lark):} Generates an Abstract Syntax Tree (AST) and maps it to the internal \texttt{QuantumCircuit} Intermediate Representation (IR). (3) \textbf{Decomposer:} Normalizes high-level gates into a target-compatible native gate set. (4) \textbf{Transpiler:} Routes qubits according to the target hardware topology using Greedy or SABRE routing algorithms. (5) \textbf{Simulators:} Executes the transpiled circuit using either the \texttt{Simulator} (exact statevector evolution) or \texttt{MPSSimulator} (tensor network contraction with SVD-based truncation). (6) \textbf{Output Layer:} Returns measurement probabilities, classical memory states, and visual plots. This modular design allows independent development, testing, and optimization of each component while maintaining a clean data flow from high-level quantum code to executable results.



% --------------------------------
% 2.5.1 Module Descriptions
% --------------------------------

\subsection{Module 1: Parser Module}

\textbf{Description}

The Parser Module serves as the entry point for quantum circuit ingestion, accepting programs in OpenQASM 2.0, OpenQASM 3.0, and JSON formats. Built on the Lark parsing toolkit, the module performs lexical analysis to tokenize input strings, syntactic analysis to construct Abstract Syntax Trees (ASTs) following formal grammar specifications, and semantic validation to ensure circuit correctness (valid qubit indices, supported gate types, proper register declarations). The parser handles advanced OpenQASM 3.0 features including classical memory operations, conditional execution, for/while loops, and timing directives. Upon successful parsing, the module generates the internal Intermediate Representation (IR) as a \texttt{QuantumCircuit} object containing structured gate sequences, qubit/classical bit registers, and metadata. Error handling provides detailed diagnostic messages with line numbers and syntax suggestions to guide users in correcting malformed circuits.

\textbf{Functional Requirements:}

\begin{itemize}
    \item FR-1.1: The system shall parse OpenQASM 2.0 circuits with support for standard gates (H, X, Y, Z, CNOT, CZ, SWAP, Toffoli) and measurement operations.
    \item FR-1.2: The system shall parse OpenQASM 3.0 circuits including classical memory declarations, conditional statements, and loop constructs.
    \item FR-1.3: The system shall validate circuit syntax and report errors with line numbers and descriptive messages.
    \item FR-1.4: The system shall generate an Intermediate Representation (IR) containing all circuit information in a standardized format.
    \item FR-1.5: The system shall support JSON circuit format for backward compatibility with simple gate list representations.
\end{itemize}



\subsection{Module 2: Transpiler Module}

\textbf{Description}

The Transpiler Module performs hardware-aware circuit compilation, transforming logical quantum circuits into physically executable programs that respect target hardware constraints. The module accepts a hardware topology specification (linear chain, 2D grid, or custom connectivity graph) and implements sophisticated qubit routing algorithms to satisfy connectivity requirements. The core functionality includes initial qubit placement using heuristic strategies, dynamic SWAP gate insertion to enable non-adjacent qubit interactions, and optional mapping restoration to preserve logical-physical qubit correspondence. The module supports two routing strategies: Greedy BFS for fast compilation with moderate optimization, and SABRE-inspired lookahead heuristics for superior circuit depth reduction (typically 30-40\% improvement). The transpiler integrates with the Decomposer to ensure all gates are expressed in the target's native gate set before routing, producing fully executable circuits ready for simulation or hardware deployment.

\textbf{Functional Requirements:}

\begin{itemize}
    \item FR-2.1: The system shall accept hardware topology specifications defining qubit connectivity constraints.
    \item FR-2.2: The system shall implement qubit routing algorithms (Greedy BFS and SABRE) to insert SWAP gates satisfying connectivity requirements.
    \item FR-2.3: The system shall minimize circuit depth and SWAP gate count through lookahead heuristics and decay mechanisms.
    \item FR-2.4: The system shall support configurable routing strategies allowing users to select optimization priorities.
    \item FR-2.5: The system shall provide optional mapping restoration to maintain logical-physical qubit correspondence after transpilation.
\end{itemize}



\subsection{Module 3: Simulator Module}

\textbf{Description}

The Simulator Module provides dual quantum state evolution engines for circuit execution and validation. The Statevector Simulator implements exact quantum mechanics using full state vector representation, supporting circuits up to 12 qubits on standard hardware through NumPy-optimized linear algebra operations. Gate application employs permutation-based indexing for efficient unitary matrix multiplication, achieving near-C performance through vectorized computation. The MPS (Matrix Product State) Simulator extends scalability to 20+ qubits for low-entanglement circuits using tensor network decomposition with SVD-based compression and configurable bond dimension truncation. Both simulators support mid-circuit measurements with probabilistic state collapse, classical memory operations (bit assignment, register manipulation), and conditional gate execution based on classical bit values. The module handles measurement sampling with configurable shot counts, probability distribution computation, and state vector extraction for analysis, providing comprehensive quantum circuit execution capabilities for algorithm development and validation.

\textbf{Functional Requirements:}

\begin{itemize}
    \item FR-3.1: The system shall simulate quantum circuits using exact statevector evolution for up to 12 qubits.
    \item FR-3.2: The system shall implement MPS simulation with SVD-based compression for low-entanglement circuits scaling to 20+ qubits.
    \item FR-3.3: The system shall support mid-circuit measurements with probabilistic state collapse and classical memory updates.
    \item FR-3.4: The system shall execute conditional gates based on classical bit values enabling hybrid quantum-classical algorithms.
    \item FR-3.5: The system shall compute measurement probability distributions and support configurable shot-based sampling.
\end{itemize}



\subsection{Module 4: Visualization Module}

\textbf{Description}

The Visualization Module generates publication-quality graphical representations of quantum circuits and simulation results using Matplotlib. The circuit diagram renderer produces professional visualizations with proper quantum gate symbols (boxes for single-qubit gates, controlled-gate notation for multi-qubit operations), horizontal qubit lines with clear labeling, measurement indicators, and classical bit wires. The module supports customizable styling including figure dimensions, color schemes, font sizes, and export formats (PNG, PDF, SVG) for integration into presentations and publications. Measurement result visualization includes probability histogram generation with configurable bar colors, axis labels, and statistical annotations. The module handles complex circuits with multiple qubits, classical registers, and conditional operations, automatically adjusting layout for readability. Visualization aids in algorithm debugging by revealing circuit structure, supports educational demonstrations by clarifying quantum operations, and enables result presentation through clear probability distributions.

\textbf{Functional Requirements:}

\begin{itemize}
    \item FR-4.1: The system shall generate circuit diagrams with proper quantum gate symbols and qubit line representations.
    \item FR-4.2: The system shall produce probability histograms visualizing measurement outcome distributions.
    \item FR-4.3: The system shall support customizable styling including colors, fonts, and figure dimensions.
    \item FR-4.4: The system shall export visualizations to multiple formats (PNG, PDF, SVG) for publication and presentation.
\end{itemize}



\subsection{Module 5: CLI and Web API Module}

\textbf{Description}

The CLI and Web API Module provides dual user interfaces for QVM access: a command-line interface for local execution and a RESTful web service for remote access. The CLI implements a comprehensive command-line tool with argument parsing for circuit file loading, transpilation configuration (routing strategy selection, topology specification), simulation parameters (seed, shot count, max operations), and output formatting options. The interface includes extensive help documentation, error reporting with actionable suggestions, and progress indicators for long-running operations. The Web API, built on FastAPI, exposes QVM functionality through HTTP endpoints with JSON request/response formatting, automatic OpenAPI documentation generation, and asynchronous request handling for concurrent circuit execution. The accompanying web dashboard provides an interactive interface for circuit composition, parameter configuration, execution monitoring, and result visualization, demonstrating the QVM's capability as a backend service for quantum computing applications and educational platforms.

\textbf{Functional Requirements:}

\begin{itemize}
    \item FR-5.1: The system shall provide a command-line interface accepting circuit files and configuration parameters.
    \item FR-5.2: The system shall implement a RESTful web API with JSON request/response formatting for remote circuit execution.
    \item FR-5.3: The system shall generate automatic API documentation using OpenAPI specifications.
    \item FR-5.4: The system shall provide an interactive web dashboard for circuit composition and result visualization.
    \item FR-5.5: The system shall support asynchronous request handling enabling concurrent circuit execution.
\end{itemize}


% --------------------------------
% 2.7 Assumptions and Dependencies
% --------------------------------

\section{Assumptions and Dependencies}

The following assumptions and dependencies guide the system design and operational constraints:

\textbf{Assumptions:}
\begin{itemize}
    \item The system assumes users have basic knowledge of quantum computing concepts (qubits, gates, superposition) and familiarity with Python programming.
    \item The system assumes quantum circuits are mathematically valid (unitary operations, proper qubit indexing) and syntactically correct according to OpenQASM specifications.
    \item The system assumes users have access to standard desktop hardware with at least 8GB RAM for simulating circuits up to 12 qubits (Statevector) or 20+ qubits (MPS for low-entanglement cases).
\end{itemize}

\textbf{Dependencies:}
\begin{itemize}
    \item The system depends on Python 3.10+ and core scientific libraries (NumPy for linear algebra, NetworkX for graph-based topology mapping, Matplotlib for visualization).
    \item The system depends on the Lark parsing library for OpenQASM 3.0 AST generation and syntax validation.
    \item The system's interoperability depends on adherence to OpenQASM 2.0 and 3.0 standards for circuit export and compatibility with external quantum platforms (e.g., IBM Quantum).
\end{itemize}

%%%%%%%%%%%%%%%%%%%%%%%%%%%%%%%%%%%%%%%%%%%%%%%%%%%%%%
\section{Chapter Summary}
%%%%%%%%%%%%%%%%%%%%%%%%%%%%%%%%%%%%%%%%%%%%%%%%%%%%%%

This chapter defined the core problem of hardware fragmentation and vendor lock-in in quantum computing, which prevents code portability across different quantum platforms. The proposed Quantum Virtual Machine (QVM) solution was presented as a middleware layer implementing a three-stage pipeline (Ingestion, Transpilation, Execution) to achieve "Write Once, Run Anywhere" functionality. Five clear business objectives were established to guide system development, focusing on hardware agnosticism, educational transparency, lightweight operation, interoperability, and automation. The system scope was defined to cover desktop-based quantum circuit simulation (10-20 qubits) using Python, with explicit boundaries excluding pulse-level control and quantum error correction. The pipeline architecture was outlined, describing the six-stage modular design from input parsing through simulation to output generation. Finally, key assumptions regarding user knowledge and hardware requirements, along with dependencies on Python libraries and OpenQASM standards, were documented to clarify operational constraints. This chapter provides the problem foundation and conceptual solution framework that will guide the detailed software design presented in the next chapter.
